\section{Introduction}

The \textit{First Proof} project~\cite{abouzaid2026}, in which eleven mathematicians contributed unpublished research problems, was designed to separate genuine reasoning from mere literature search. Its second batch, evaluated by human judges~\cite{firstproof2b}, documented a concrete failure mode of language models applied to mathematics: systems performed better when the problem was structurally similar to an already published result, and in several cases solved it by translating a prior proof from the literature, reusing even its terminology and notation line by line, without citing the source, to the point that a human reviewer would have flagged it as plagiarism. What is missing, then, is not the ability to generate correct proofs, but a systematic method to assess whether a proof is genuinely new. Models can generate statements that compile, that are correct, and yet are not new.

The problem is twofold. On the one hand, current mathematical AI systems (Axiom, Harmonic, DeepSeek-Prover) verify correctness but not novelty. On the other, the naive criterion that equates novelty with absence from Mathlib~\cite{kasaura2025} fails in two directions: it marks as new trivial theorems that any automated tactic from 2026 closes without requiring a genuine mathematical idea, and it misses results that are not formalized but are published in the informal literature.

Tao~\cite{tao2025}, in his notes on \textit{Machine-Assisted Proof}, frames the problem in terms of scale: formal verification is human-labor intensive, to the point that today it is not feasible to formalize in real time a significant fraction of research articles, while automated assistance promises mathematical exploration at scales currently unattainable. Tao directs that argument to the control of \emph{correctness} and to the scarcity of reviewers willing to verify long proofs line by line. The same argument, in our view, transfers to the control of \emph{novelty}: if theorem generation accelerates, manually evaluating whether each result is truly new becomes as infeasible as that line-by-line verification, and the one proposing the theorems becomes a model.

The urgency of the problem was illustrated while this work was being written. In July 2026, Levent Alp\"oge, a mathematician at Anthropic, announced an explicit counterexample to the Jacobian conjecture, open since Keller formulated it in 1939, found with the assistance of a language model and so brief that its verification is elementary: a symbolic calculation of seconds, independently formalized in Lean and Isabelle within hours against a pre-registered statement of the conjecture~\cite{alpoge2026}. The counterexample extends immediately to all dimensions $\geq 3$, and within days families of higher-dimensional counterexamples were circulating; the episode is also part of a broader wave of AI-assisted refutations of other open conjectures. Such an episode produces in concentrated form the material that motivates this work: results descending from the same discovery, obtained in parallel by different authors, where the question of whether a statement was already established by another is genuine and costly to answer manually. It is furthermore the regime in which the coverage limitation we document in section~\ref{ssec:cobertura} does not apply: the potentially duplicated works are contemporary and of immediate public access.

This article presents AViD Journal (Automated Verification in Demonstrations), a system that addresses that question operationally: it receives a LaTeX article, formalizes its statements in Lean~4, and issues a novelty verdict through a three-dimensional decision tree.

\subsection{Contributions}

\begin{enumerate}
\item \textbf{A complete novelty verification pipeline.} The system traverses the entire path, from LaTeX source to verdict: parsing of mathematical environments, formalization in Lean~4 with multi-model abstraction, and an eight-verdict decision tree over three dimensions: prior existence in a formal corpus (Mathlib, via Leandex) and an informal one (TheoremSearch and Matlas, with temporal filter and LLM judge), non-triviality via Lean's automatic tactics, and structural distance between proofs as Jaccard distance over premise sets.

\item \textbf{A benchmark of papers withdrawn for duplication.} A dataset of papers withdrawn from arXiv whose withdrawal comment declares duplication of prior results, with controls matched by category and year, and the analysis of the retrievability of their duplicators.

\item \textbf{A map of the failure modes of the approach.} The central contribution of the evaluation is not a performance metric but the characterization of three walls: (a) that a Lean file compiles does not guarantee that the formalization is faithful to the original statement; (b) the results that duplication rediscovers tend to predate the preprint era, and are therefore invisible to available corpora; and (c) arXiv removes the LaTeX source upon withdrawing a paper, which limits the reproducibility of any benchmark built upon them.
\end{enumerate}

\subsection{Structure of the article}

Section~\ref{sec:related} situates this work with respect to the research fronts that intersect in novelty verification: statement search infrastructure, novelty filters in conjecture generation, proof identity, and the distinction between correctness and novelty. Section~\ref{sec:pipeline} describes the pipeline: parser, formalization, reference corpora, the three dimensions, and the verdict tree. Section~\ref{sec:instrumentos} validates each instrument separately. Section~\ref{sec:experimentos} reports the evaluation on withdrawn papers and characterizes the failure modes found. Section~\ref{sec:limitaciones} enumerates the limitations of the framework, the implementation, and the experimental design. Section~\ref{sec:conclusion} closes with what has been built and what remains pending.